\part{理论延伸：同一个算法，两种剧本}

\section{学历上层：被聪明和脚手架掩盖的 ADHD}

先讲一个反直觉的事实：\term{学历上层反而是 ADHD 最隐蔽的地方}。

\subsection{画像}

一路名校，硕士博士，体面的职业。小时候「看着做的事挺多、成绩还行」，从没被任何人怀疑过。可能直到二十多岁甚至三十多岁，在某次崩溃、某段关系破裂，或某个视频下面，才第一次认真怀疑：我是不是 ADHD？

为什么这么晚才发现？因为他们的算法缺陷，被两样东西盖住了：\term{聪明}，和\term{脚手架}。

\subsection{他们是怎么「混」过来的}

\begin{itemize}
  \item \term{智力补丁}：低要求阶段，靠聪明直接覆盖执行缺陷。小学初中内容简单，被动注意力够用，不用主动注意力也能学个七七八八。
  \item \term{肾上腺素冲刺}：靠死线启动。D 前三天是天才，D 前三个月是废物——这不是夸张，这是常态。紧迫感是唯一有效的启动器，于是每件事都被他们亲手拖到变成紧急事。
  \item \term{超专注红利}：偶尔对某件事上头，一次超专注顶普通人几个月的线性积累。以广促深在长周期里确实成立，所以他们能「平时不学，考前封神」。
  \item \term{环境脚手架}：好学校、好家庭、好导师，处处有人替他们把结构架好。课表有人排，作业有人收，死线有人催。他们的执行功能缺憾，被环境无声地覆盖了。
\end{itemize}

\subsection{习惯观察}

学历上层的 ADHD，日常往往是这样一副矛盾的样子：

\begin{checklist}
  \item \textbf{死线前爆发}；能拖则拖；对「无趣但有长期回报」的任务完全无感；
  \item \textbf{生活层面的功能崩塌}：堆积的快递、没拆的包装、未回的微信、未缴的账单、乱到没法收拾的房间。工作体面，生活失能；
  \item \textbf{自我药物化}：咖啡因、酒精、短视频、游戏、报复性熬夜——靠刺激续命；
  \item \textbf{完美主义面具}：越聪明越不敢交「不完美」的东西，于是更拖。别人催他，他说「还没准备好」，其实是在等那个永远等不到的「准备好了」。
\end{checklist}

\subsection{成因：环境把执行缺陷盖住了}

这里要说透一个扎心的事实：\term{他们不是没问题，是问题被环境买断了}。学校的课表、导师的催逼、家庭的后勤，都是外部脚手架——是环境在替他们的前额叶上班。这解释了为什么很多高功能 ADHD 一脱离校园（进入自由度极高的科研、自由职业，或独居生活）就立刻崩掉：脚手架撤了，缺陷第一次裸奔。

\begin{insight}
学历上层的核心矛盾：\term{能力天花板很高，燃料系统却在慢性透支}。他们是靠肾上腺素而不是靠系统运转的——每一次冲刺都是透支，多巴胺阈值越抬越高，于是越来越「撑不住」。很多人的「中年危机式的觉醒」，本质就是这个引擎烧到冒烟了。
\end{insight}

\subsection{代价：燃尽、冒充者、与越成功越空心}

\begin{itemize}
  \item \term{燃尽循环}：冲刺—透支—恢复—再冲刺，周而复始，直到某次彻底转不动。高功能 ADHD 最常见的就医触发点，不是「做不好」，而是「突然做不动了」。
  \item \term{冒充者综合征}：总觉得自己是「靠运气和最后一晚混过来的」，随时可能被拆穿。这个想法不完全是错的——他们确实在靠最后一晚——但代价是不敢停、不敢露怯、不敢求助。
  \item \term{越成功越空心}：成绩、职位带来的多巴胺在边际递减。小学考第一还有奖励，读完博士发完顶刊，奖励系统彻底失灵——「然后呢？」「这有啥用？」这正是第三章讲的动机悖论在顶端的显影。
  \item \term{恐惧失败导致不敢开始}：因为太聪明，太知道「这件事如果认真做能做成什么样」，反而被那个标准吓住，于是宁可拖着。
\end{itemize}

\subsection{一句话总结这一章}

学历上层不是「没事」，是\term{一直有事，只是还没崩}。他们的自救起点，不是「学会更努力」，而是承认自己需要脚手架——哪怕外表看起来一切正常。
